%XXXXXXXXXXXXXXXXXXXXXXXX-----Kapitel-----XXXXXXXXXXXXXXXXXXXXXXXXXXX
%XXXXXXXXXXXXXXXXXXXXXXXXXXXXXXXXXXXXXXXXXXXXXXXXXXXXXXXXXXXXXXXXXXXX
%XXXXXXXXXXXXXXXXXXXXXXXXXXXXXXXXXXXXXXXXXXXXXXXXXXXXXXXXXXXXXXXXXXXX
%Kapitel Vorlage für Bilder, Formeln, etc......
\chapter{Vorlagen für LaTeX}
\label{Kapitel_Vorlage}
%XXXXXXXXXXXXXXXXXXXXXXXXXXXXXXXXXXXXXXXXXXXXXXXXXXXXXXXXXXXXXXXXXXXX
%XXXXXXXXXXXXXXXXXXXXXXXXXXXXXXXXXXXXXXXXXXXXXXXXXXXXXXXXXXXXXXXXXXXX
%XXXXXXXXXXXXXXXXXXXXXXXXXXXXXXXXXXXXXXXXXXXXXXXXXXXXXXXXXXXXXXXXXXXX


%XXXXXXXXXXXXXXXXXXXXXXXX-----Sektion-----XXXXXXXXXXXXXXXXXXXXXXXXXXX
%XXXXXXXXXXXXXXXXXXXXXXXXXXXXXXXXXXXXXXXXXXXXXXXXXXXXXXXXXXXXXXXXXXXX
\section{Sektion Vorlage}
\label{Sektion_Vorlage}
%XXXXXXXXXXXXXXXXXXXXXXXXXXXXXXXXXXXXXXXXXXXXXXXXXXXXXXXXXXXXXXXXXXXX
%XXXXXXXXXXXXXXXXXXXXXXXXXXXXXXXXXXXXXXXXXXXXXXXXXXXXXXXXXXXXXXXXXXXX

%XXXXXXXXXXXXXXXXXXXXXXXX-----Subsektion--XXXXXXXXXXXXXXXXXXXXXXXXXXX
\subsection{Subsektion Vorlage}
\label{Subsektion_Vorlagen}
%XXXXXXXXXXXXXXXXXXXXXXXXXXXXXXXXXXXXXXXXXXXXXXXXXXXXXXXXXXXXXXXXXXXX


%XXXXXXXXXXXXXXXXXXXXXXXX-----MiniSek--------XXXXXXXXXXXXXXXXXXXXXXXX
\minisec{MiniSek} \label{MiniSek_Vorlagen}
%XXXXXXXXXXXXXXXXXXXXXXXXXXXXXXXXXXXXXXXXXXXXXXXXXXXXXXXXXXXXXXXXXXXX

%XXXXXXXXXXXXXXXXXXXXXXXX-----Paragraf-------XXXXXXXXXXXXXXXXXXXXXXXX
\paragraph{Paragraf Vorlage} \label{Paragraf_}
%XXXXXXXXXXXXXXXXXXXXXXXXXXXXXXXXXXXXXXXXXXXXXXXXXXXXXXXXXXXXXXXXXXXX


%-----------------------------Gleichung------------------------------
\begin{equation}
 \label{Gleichung_1}
\end{equation}
%--------------------------------------------------------------------


%-----------------------------mehrere Gleichungen--------------------
\begin{eqnarray}
\label{Gleichung_2}
 \nonumber
 &=&
 \\
 &=&
 \\
\end{eqnarray}
%--------------------------------------------------------------------


%-----------------------------Aufzählung-----------------------------
\begin{Aufzaehlung}{$\bullet$}
\item
blabla
\item
blabla
\end{Aufzaehlung}
%--------------------------------------------------------------------



%Ein Bild kann man so "`einbinden"': (Der Ausdruck "`clip="'
%verhindert, dass das Bild über seine
%Boundingbox überläuft, das hei{\ss}t es ist auf keinen Fall grö{\ss}er!)
%


%///////////////////////////// Bild /////////////////////////////////
\begin{figure}[h!]\begin{center}
 \includegraphics[width=\BildBreiteGross, clip=]{Calvin_und_Hobbes2}
 %[width=15cm,height=8cm,scale=0.5,angle=45]
 \caption{Calvin und Hobbes \cite{Watterson2013}}
 %\label{Bild_}
\end{center}\end{figure}
%////////////////////////////////////////////////////////////////////

%///////////////////////////// 2 Bilder nebeneinander ///////////////
\begin{figure}[h!]\centering
 \subfloat[Bildbeschreibung subfigure 1]
    {\label{Bild_subfigure_1}
    \includegraphics[width=\BildSubfigureBreite]{Calvin_und_Hobbes2}}
 \subfloat[Bildbeschreibung subfigure 2]
    {\label{Bild_subfigure_2}
    \includegraphics[width=\BildSubfigureBreite]{Calvin_und_Hobbes2}}
 \caption{Hauptbildunterschrift}
\end{figure}
%////////////////////////////////////////////////////////////////////

\begin{figure}[h!]\centering
\captionsetup[subfloat]{labelformat=empty} % (a), (b) unterdrücken
 \subfloat[]
    {\label{fig:Bild_subfigure_1}
    \includegraphics[width=0.45\textwidth]{Calvin_und_Hobbes2}}
    \hspace{0.5 cm}
 \subfloat[]
    {\label{fig:Bild_subfigure_2}
    \includegraphics[width=0.45\textwidth]{Calvin_und_Hobbes2}}
 \caption{Zählpfeile und Vorzeichen des Stromes}
 \captionsetup[subfloat]{labelformat=parens}
\end{figure}

% Beschrifttung unterdrückt

%-----------------------------Tabelle--------------------------------
\begin{table}[h!]
  \centering
  \renewcommand{\arraystretch}{1.5}
  \begin{tabular}{|p{5.3cm}|p{10cm}|} \hline
    t&t
    \\\hline\hline
    t&t
    \\\hline
    \end{tabular}
  \caption{Versuchstabelle}
  \label{Tabelle_Versuch}
\end{table}
%--------------------------------------------------------------------




%-----------------------------Aufzählung-----------------------------
\begin{list}{\bullet}{
 \setlength{\leftmargin}{4cm}
 \setlength{\labelsep}{1cm}
 \setlength{\labelwidth}{4cm}}

\item[Punkt 1] bla bla bla ldfkjasdhfjlk hjkadhfas jk
dkashfjkas hjkdf hajks ajd jkf

\item[Punkt 2]

\item[Punkt 3]

\end{list}
%--------------------------------------------------------------------





Rotierter Text (erst in ps oder pdf sichtbar):

\begin{rotate}{90}
Rotiert?
\end{rotate}

\vspace{2cm}

\begin{rotate}{90}
$H + 4 -\frac{3}{4}$
\end{rotate}




$_K\vec{r}_{C2} = [l_{C2} \quad 0 \quad 0 ]^T$

 $_I\vec{e}_1$

Gänsefüsschen   "`Text"'

Einheiten:

$[\frac{kgm^2}{s}]$ für die Dämpfung

$[\frac{rad}{s}]$ für Winkelgeschwindigkeiten

$[\frac{kgm^2rad}{s^2}]$ für die Reibung



$_I\vec{e}_1$

$\dot{q}_2x$


Link auf eine Homepage
\href{http://de.wikipedia.org/wiki/Welt}{hier}



Matrizen kann man als "`Tabellen"' aufschreiben:
%--------------------------------------------------------------------
 %\begin{displaymath} {^{KI}\mf{A}=}
%\label{Matrix_A_K_I_Beispiel} \left(
%\begin{array}{ccc}

%\cos{q_2}\cos{\beta}+\sin{q_2}\sin{\alpha}\sin{\beta}
%&\sin{q_2}\cos{\alpha}
%&-\cos{q_2}\sin{\beta}+\sin{q_2}\sin{\alpha}\cos{\beta}\\

%-\sin{q_2}\cos{\beta} + \cos{q_2}\sin{\alpha}\sin{\beta}
%&\cos{q_2}\cos{\alpha}
%&\sin{q_2}\sin{\beta}+\cos{q_2}\sin{\alpha}\cos{\beta}\\

%\cos{\alpha}\cos{\beta} &-\sin{\alpha}
%&\cos{\alpha}\cos{\beta}

%\end{array}
%\right)
%\end{displaymath}\\
%--------------------------------------------------------------------



Matrizen kann man besser so aufschreiben:
%--------------------------------------------------------------------
%\begin{equation}
%^{KI}\mf{A}=
%\begin{pmatrix}
%
%\cos{\beta}        &0       &-\sin{\beta}\\
%
%0                  &1       &0\\
%
%\sin{\beta}        &0       &\cos{\beta}
%
%\end{pmatrix}
%\end{equation}
%--------------------------------------------------------------------


Für Vektoren
%--------------------------------------------------------------------
\begin{equation}\label{Matrix_A_K_alpha_Beispiel}
\boxed{ _K\vec{r}_{C2}=
\begin{pmatrix}
ls2\\0\\0
\end{pmatrix}}
\end{equation}
%--------------------------------------------------------------------


Aufzählung selbst gestalten:
\begin{list}{$\bullet$}{\setlength
        {\leftmargin}{0.5cm}}
        \item Zusatzlenkwinkel durch das AFS-System
        \item Stellung eines Zusatzmoments
\end{list}




%XXXXXXXXXXXXXXXXXXXXXXXXXXXXXXXXXXXXXXXXXXXXXXXXXXXXXXXXXXXXXXXXXXXX
%XXXXXXXXXXXXXXXXXXXXXXXXXXXXXXXXXXXXXXXXXXXXXXXXXXXXXXXXXXXXXXXXXXXX
\section{Sektion Versuche}
\label{Sektion_Versuche}
%XXXXXXXXXXXXXXXXXXXXXXXXXXXXXXXXXXXXXXXXXXXXXXXXXXXXXXXXXXXXXXXXXXXX
%XXXXXXXXXXXXXXXXXXXXXXXXXXXXXXXXXXXXXXXXXXXXXXXXXXXXXXXXXXXXXXXXXXXX
Test für Formelzeichen: Im Text funktioniert diese mbox gut:
\mbox{$x = y$}.
%-----------------------------Gleichung------------------------------
\begin{equation}
 Sogar in Formeln kann man die mbox mit integrierter Formel
 verwenden:
 x = y, \mbox{$x = y$}
\end{equation}
%--------------------------------------------------------------------


Partialbruchzerlegung:
%--------------------------------------------------------------------
\begin{eqnarray*}
 f(t) = \frac{As + B}{s^3(J s^2 + d s + c)} &=&
 -{\frac {cJB+dcA-{d}^{2}B}{{c}^{3}s}}
 \\
 &&
 +{\frac {cA-dB}{{c}^{2}{s}^{2}}}
 +{\frac {B}{c{s}^{3}}}
 \\
 &&
 +{\frac
 {-{d}^{3}B-J{c}^{2}A+2\,cJdB+{d}^{2}cA+s{J}^{2}cB+sJdcA
 -sJ{d}^{2}B}{{c}^{3}\left
 (J{s}^{2}+ds+c\right )}}
\end{eqnarray*}
%--------------------------------------------------------------------
